\subsection{Multiple Radical Terms} 
We noted before that when an expression has one term with a radical and one term without, we can't combine those terms. When an expression has multiple radical terms, we can simplify by combining ``like radicals.'' It is useful to first fully simplify each radical term.
\begin{Example}
Simplify $\sqrt{90}-3\sqrt{20}+2\sqrt{\frac{18}{5}}+4\sqrt{80}$.
\begin{lpic}{}{8}
\foreach \y/\exp in {1/90,2/20,5/80}
	\regtextright{2.5}{-\y}{$\sqrt{\exp}={}$};
\fractextright{2.5}{-3}{$\sqrt{\dfrac{18}{5}}={}$};
\regtextleft{0.25}{-6}{$\sqrt{90}-3\sqrt{20}+2\sqrt{\frac{18}{5}}+4\sqrt{80}={}$};
\end{lpic}
We can check that this worked by taking a decimal approximation of the original expression and the simplified expression.
\begin{lpic}{}{2}
\regtextleft{0.25}{-1}{$\sqrt{90}-3\sqrt{20}+2\sqrt{\frac{18}{5}}+4\sqrt{80}\approx{}$};
\end{lpic}
\end{Example}
\noi When terms include both variables and radicals, both must be exactly the same to combine the terms.
\begin{Example}
Simplify $xy^2\sqrt{20x^5y}+3x^2y\sqrt{45x^3y^3}$. Assume all variables represent positive real numbers.
\begin{lpic}{}{5}
\foreach \y/\exp in {1/20x^5y,2/45x^3y^3}
	\regtextright{3.5}{-\y}{$\sqrt{\exp}={}$};
\end{lpic}
\end{Example}